% ==============================================================================
% 05_forensics_fondements.tex — UE SIS589
% ==============================================================================
\chapter{Fondements de l'Investigation Numérique}
\label{chap:forensics-fondements}

\epigraph{\itshape « Every contact leaves a trace. »}
         {--- Edmond Locard, Principe de transfert}

\begin{boxobjectifs}
\begin{enumerate}
  \item Définir l'investigation numérique et ses objectifs (technique et juridique).
  \item Appliquer l'ordre de volatilité RFC~3227 à un scénario d'incident.
  \item Acquérir une image disque bit à bit avec vérification d'intégrité.
  \item Établir et maintenir une chaîne de custody conforme à ISO~27037.
  \item Distinguer les types de données forensiques~: allouées, supprimées,
        espace libre, slack space, métadonnées.
  \item Utiliser les outils d'acquisition de base~: \texttt{dd}, FTK~Imager,
        \texttt{dcfldd}.
\end{enumerate}
\end{boxobjectifs}

% ==============================================================================
\section{Définition et objectifs de l'investigation numérique}
\label{sec:forensics-def}
% ==============================================================================

\begin{boxdef}
\textbf{Investigation numérique (\textit{Digital Forensics})~:} Science
consistant à identifier, préserver, analyser et présenter des preuves
numériques de manière méthodique, reproductible et légalement recevable.
Elle répond à deux objectifs complémentaires~:
\begin{itemize}
  \item \textbf{Technique~:} comprendre comment l'attaque s'est déroulée,
        identifier la vulnérabilité exploitée, reconstruire la timeline,
        permettre la remédiation.
  \item \textbf{Juridique~:} constituer un dossier de preuves recevables
        pour une action pénale, civile ou disciplinaire.
\end{itemize}
\end{boxdef}

L'investigation numérique se distingue d'une simple analyse technique par
son \textbf{impératif d'intégrité}~: chaque action est documentée, chaque
preuve est scellée, chaque outil est qualifié. Un examen bâclé peut détruire
des preuves légalement irremplaçables.

\subsection{Domaines de l'investigation numérique}

\begin{table}[H]
\centering\small
\caption{Domaines de l'investigation numérique}
\label{tab:forensics-domaines}
\rowcolors{2}{TableauPair}{TableauImpair}
\begin{tabular}{p{3.5cm}p{5.5cm}p{4.5cm}}
\toprule
\tableauHeader{Domaine} & \tableauHeader{Objet} &
\tableauHeader{Outils de référence} \\
\midrule
Disque / Système de fichiers &
Artefacts OS, fichiers supprimés, timeline &
Autopsy, FTK, The Sleuth Kit \\
Mémoire vive (RAM) &
Processus, connexions réseau, clés en mémoire, malwares injectés &
Volatility~3, Rekall \\
Réseau &
Capture et analyse de trafic, reconstitution sessions &
Wireshark, NetworkMiner, Zeek \\
Mobile &
Applications, SMS, localisation, données cloud &
Cellebrite, Magnet AXIOM, Oxygen \\
Cloud &
Logs cloud, API, IAM, buckets S3 &
CloudTrail, Azure Monitor, GCP Logging \\
Malware (RE) &
Désassemblage, débogage, comportement sandbox &
Ghidra, IDA Pro, Cuckoo Sandbox \\
\bottomrule
\end{tabular}
\end{table}

% ==============================================================================
\section{Le principe fondamental~: ne rien altérer}
\label{sec:forensics-alteration}
% ==============================================================================

\begin{boiteTheoreme}{Principe d'intégrité forensique}
Toute interaction avec un système susceptible de contenir des preuves doit
être conçue de sorte à ne laisser aucune trace modifiant l'état de ces preuves.
Si une modification est techniquement inévitable (mise sous tension, accès
mémoire), elle doit être minimisée, justifiée et intégralement documentée.
\end{boiteTheoreme}

En pratique, ce principe implique~:
\begin{itemize}
  \item Utiliser des \termecle{bloqueurs en écriture} (\textit{write blockers})
        matériels lors de l'accès aux disques.
  \item Travailler systématiquement sur des \textbf{copies forensiques}
        (images), jamais sur l'original.
  \item Calculer et vérifier les hashes d'intégrité avant et après chaque
        étape.
  \item Documenter chaque action dans un registre horodaté.
\end{itemize}

% ==============================================================================
\section{Ordre de volatilité et stratégie d'acquisition}
\label{sec:volatilite}
% ==============================================================================

\subsection{Ordre de volatilité (RFC~3227)}

L'ordre de collecte est dicté par la durée de vie des données~: les données
les plus éphémères sont collectées en premier.

\begin{table}[H]
\centering\small
\caption{Ordre de volatilité avec protocole d'acquisition}
\label{tab:volatilite-protocole}
\rowcolors{2}{TableauPair}{TableauImpair}
\begin{tabular}{p{0.7cm}p{4cm}p{4cm}p{4.5cm}}
\toprule
\tableauHeader{\#} & \tableauHeader{Donnée} &
\tableauHeader{Outil d'acquisition} & \tableauHeader{Remarque} \\
\midrule
1 & Registres CPU, cache & Non capturable sans matériel spécialisé & Accepter la perte \\
2 & Mémoire vive (RAM) & \texttt{WinPmem}, \texttt{avml}, \texttt{LiME} &
  Avant tout autre action sur le système \\
3 & Tables réseau (ARP, routage, connexions) &
  \texttt{netstat -anp}, \texttt{arp -a}, \texttt{ss -tnp} &
  Rapide, quelques lignes de commandes \\
4 & Processus en cours & \texttt{ps auxf}, \texttt{lsof}, \texttt{tasklist} &
  Avec PID, PPID, arguments \\
5 & Disque dur / SSD & \texttt{dd}, \texttt{dcfldd}, FTK~Imager &
  Image bit à bit avec bloqueur en écriture \\
6 & Journaux distants & Export SIEM, serveur Syslog distant & Déjà collectés en temps réel \\
7 & Archives / backups & Copie simple & Persistant \\
\bottomrule
\end{tabular}
\end{table}

\subsection{Décision~: système éteint ou allumé~?}

\begin{description}
  \item[\textbf{Système allumé (\textit{live acquisition})~:}]
        Permet de capturer la mémoire vive, les connexions réseau actives,
        les processus, les fichiers ouverts, les clés de chiffrement en mémoire.
        \textbf{Recommandé si le système est critique ou si le malware est actif.}
        Risque~: l'attaquant peut avoir un programme de destruction de preuves
        déclenché par une activité anormale.
  \item[\textbf{Système éteint (\textit{dead acquisition})~:}]
        Garantit l'intégrité du disque~; perd la mémoire vive et l'état
        réseau. \textbf{Recommandé pour l'acquisition du disque.}
        Protocole~: débrancher l'alimentation sans arrêt logiciel (pour
        éviter l'exécution des scripts d'arrêt qui pourraient effacer).
\end{description}

% ==============================================================================
\section{Acquisition d'une image disque}
\label{sec:acquisition}
% ==============================================================================

\subsection{Principes}

L'image forensique est une copie \textbf{bit à bit} de l'intégralité du
support de stockage~: secteurs alloués, secteurs libres, espace de swap,
secteur de démarrage, données supprimées. Elle est différente d'une copie
de fichiers (qui ignore les données supprimées et le slack space).

\subsection{Formats d'image}

\begin{description}
  \item[\textbf{RAW (\texttt{.dd}, \texttt{.img})~:}]
        Format brut~: copie bit à bit sans métadonnées. Simple, universel.
        Taille égale à la taille du disque source.
  \item[\textbf{E01 (EnCase Evidence File)~:}]
        Format propriétaire (Guidance Software / OpenEvidence).
        Inclut les métadonnées d'acquisition, les hashes MD5/SHA-1,
        et la compression. Standard de facto en environnement judiciaire.
  \item[\textbf{AFF4 (\textit{Advanced Forensic Format})~:}]
        Format ouvert, supporte la compression, le chiffrement, les
        annotations. Utilisé par Velociraptor et certains outils DFIR modernes.
\end{description}

\subsection{Acquisition avec \texttt{dd} et \texttt{dcfldd}}

\begin{lstlisting}[style=StyleTerminal,
  caption={Acquisition forensique Linux --- dd et dcfldd},
  label={lst:dd-acquisition}]
# ─── Identifier le disque cible (SANS le monter) ────────────────────────────
lsblk
fdisk -l
# Résultat : /dev/sdb = disque cible (NE PAS MONTER)

# ─── dd classique : copie bit à bit ─────────────────────────────────────────
dd if=/dev/sdb \
   of=/mnt/evidence/srv-web-01-sda.dd \
   bs=64K \
   conv=noerror,sync \
   status=progress

# Options :
# if=   : source (disque cible, JAMAIS la machine live sans write blocker)
# of=   : destination (disque externe ou réseau)
# bs=   : taille de bloc (64K = optimal SSD)
# conv=noerror,sync : continuer même si erreur de lecture (bad sectors)
# status=progress   : afficher l'avancement

# ─── dcfldd : version forensique de dd (hash intégré) ───────────────────────
dcfldd if=/dev/sdb \
       of=/mnt/evidence/srv-web-01-sda.dd \
       hash=sha256 \
       hashlog=/mnt/evidence/srv-web-01-sda.hash \
       bs=64K \
       conv=noerror,sync \
       sizeprobe=if

# ─── Vérification d'intégrité après acquisition ──────────────────────────────
sha256sum /dev/sdb > source.sha256
sha256sum /mnt/evidence/srv-web-01-sda.dd > image.sha256
diff source.sha256 image.sha256   # DOIT être identique
\end{lstlisting}

\subsection{Acquisition avec FTK Imager (Windows)}

FTK Imager (AccessData, gratuit) est l'outil standard en environnement
Windows pour~:
\begin{itemize}
  \item Créer des images E01 ou RAW avec calcul automatique MD5/SHA-1.
  \item Capturer la mémoire vive (\textit{memory dump}).
  \item Prévisualiser le contenu d'une image sans la monter.
  \item Exporter des fichiers individuels depuis une image.
\end{itemize}

\begin{boxcustody}
\textbf{PV d'acquisition --- Protocole minimal}

\begin{itemize}
  \item Date/heure (UTC)~: \timestamp{2026-04-12 04:22:00}
  \item Enquêteur~: MINKA T.E. --- Analyste N2 SOC
  \item Système source~: srv-web-01.lizbiz.local (IP~: 10.0.1.10)
  \item Matériel source~: SSD SATA 500~Go, SN~: WD-20260101ABCD
  \item Bloqueur en écriture~: Tableau UltraBlock SATA/IDE, SN~: TB-0042
  \item Outil~: dcfldd v1.3.4-1 sur Ubuntu~22.04 (SIFT Workstation)
  \item Destination~: \texttt{srv-web-01-sda.dd} sur disque externe
        chiffré VeraCrypt, SN~: EXT-2026-04
  \item \hash{sha256}{3a7d4f8bc2e910...} (source)
  \item \hash{sha256}{3a7d4f8bc2e910...} (image) --- \textbf{IDENTIQUES}
  \item Conservation~: coffre-fort salle ~133, scellé n°~2026-04-001
\end{itemize}
\end{boxcustody}

% ==============================================================================
\section{Acquisition de la mémoire vive}
\label{sec:acq-ram}
% ==============================================================================

\begin{boxalert}
\textbf{La mémoire vive est la preuve la plus précieuse et la plus éphémère.}
Elle contient~: processus en cours d'exécution, connexions réseau actives,
clés de chiffrement non protégées, artefacts de malwares injectés,
mots de passe en clair pour certains services. Elle disparaît à l'extinction
du système. \textbf{Priorité n°~1 en \textit{live acquisition}.}
\end{boxalert}

\begin{lstlisting}[style=StyleTerminal,
  caption={Acquisition mémoire vive --- Linux (avml) et Windows (WinPmem)},
  label={lst:ram-acquisition}]
# ─── Linux : avml (Microsoft, open-source, recommandé) ──────────────────────
# Télécharger avml sur la machine d'analyse, copier sur clé USB
# Exécuter DEPUIS la clé (ne pas installer sur la machine cible)
sudo ./avml /mnt/usb/ram-dump-srv-web-01-$(date +%Y%m%d-%H%M%S).lime

# Vérifier l'intégrité
sha256sum /mnt/usb/ram-dump-srv-web-01-*.lime

# ─── Linux : LiME (Loadable Kernel Module) ───────────────────────────────────
# Exige la compilation du module pour le kernel exact de la cible
# À faire sur machine identique, pas sur la cible
insmod lime-$(uname -r).ko "path=/mnt/usb/ram.lime format=lime"

# ─── Windows : WinPmem ───────────────────────────────────────────────────────
# Depuis une invite admin, sur un support externe
winpmem_mini_x64.exe ram-dump-srv-web-01.raw
# Hash immédiatement après :
certutil -hashfile ram-dump-srv-web-01.raw SHA256

# ─── FTK Imager (Windows, GUI) ───────────────────────────────────────────────
# File > Capture Memory > Destination : E:\ (support externe)
# ✓ Include pagefile (inclure le fichier de pagination)
# → Produit un fichier .mem avec hash MD5/SHA-1
\end{lstlisting}

% ==============================================================================
\section{Structures de données forensiques}
\label{sec:forensics-data}
% ==============================================================================

\begin{table}[H]
\centering\small
\caption{Types de données forensiques sur disque}
\label{tab:forensics-data}
\rowcolors{2}{TableauPair}{TableauImpair}
\begin{tabular}{p{3.5cm}p{5.5cm}p{4.5cm}}
\toprule
\tableauHeader{Type} & \tableauHeader{Description} &
\tableauHeader{Intérêt forensique} \\
\midrule
Données allouées &
Fichiers et répertoires actuellement présents dans le système de fichiers &
Accès direct~; état courant du système \\
Données supprimées &
Fichiers supprimés dont les secteurs ne sont pas encore réalloués &
Récupération possible~; cache de fichiers effacés \\
Espace non alloué (\textit{unallocated}) &
Secteurs libres ou anciennement occupés &
Récupération de fichiers effacés, carving \\
\textit{Slack space} &
Espace résiduel entre la fin d'un fichier et la fin de son cluster &
Peut contenir des fragments de données antérieures \\
Métadonnées &
Timestamps, permissions, propriétaire, taille, inode/MFT &
Timeline, attribution, vérification d'altération \\
Fichiers de swap &
Mémoire virtuelle déversée sur disque (pagefile.sys, swap) &
Fragments de processus, données en mémoire \\
\bottomrule
\end{tabular}
\end{table}

\begin{boxdef}
\textbf{File Carving~:} Technique de récupération de fichiers supprimés
reposant sur la recherche de signatures de fichiers connues (\textit{magic bytes})
dans l'espace non alloué, sans s'appuyer sur les métadonnées du système de
fichiers (qui peuvent être effacées). L'outil \texttt{Scalpel} et Autopsy
implémentent cette technique.
\end{boxdef}

% ==============================================================================
\section{Environnement d'analyse forensique}
\label{sec:forensics-env}
% ==============================================================================

\begin{boxdef}
\textbf{SIFT Workstation (SANS Institute)~:} Distribution Linux préconfigurée
regroupant l'ensemble des outils forensiques open-source~: Autopsy, Volatility,
log2timeline, Plaso, The Sleuth Kit, NetworkMiner, bulk\_extractor, etc.
Disponible en VM VMware ou OVA. Standard de formation GIAC/SANS.
\end{boxdef}

Règles d'or de l'environnement d'analyse~:
\begin{itemize}
  \item \textbf{Isolement réseau~:} la station d'analyse est isolée
        d'Internet pendant l'analyse (risque de beacon-back).
  \item \textbf{Snapshot avant chaque analyse~:} permettre le retour
        arrière si une erreur est commise.
  \item \textbf{Journalisation des commandes~:} tout ce qui est tapé
        est enregistré (\texttt{script session.log}).
  \item \textbf{Ne jamais monter une image en lecture-écriture.}
\end{itemize}

\begin{lstlisting}[style=StyleTerminal,
  caption={Montage d'une image disque en lecture seule (SIFT)},
  label={lst:mount-image}]
# Monter une image RAW en lecture seule avec loopback
sudo losetup -r -f srv-web-01-sda.dd
losetup -a   # vérifier le nom assigné, ex: /dev/loop0

# Identifier les partitions dans l'image
sudo fdisk -l /dev/loop0
mmls srv-web-01-sda.dd   # The Sleuth Kit

# Monter la partition (offset en secteurs × 512 = bytes)
# Exemple : partition Linux ext4 commence au secteur 2048
sudo mount -o ro,loop,offset=$((2048*512)) \
           srv-web-01-sda.dd /mnt/forensics/

# Ou directement avec Autopsy (GUI) : Add Evidence Item > Image File
# Autopsy calcule automatiquement les offsets et monte en lecture seule.

# Toujours vérifier le hash AVANT d'analyser
sha256sum srv-web-01-sda.dd  # Doit correspondre au PV d'acquisition
\end{lstlisting}

% ==============================================================================
\section*{Résumé}
% ==============================================================================

\begin{boxresume}
\begin{itemize}
  \item L'investigation numérique a deux objectifs~: technique
        (comprendre l'attaque) et juridique (produire des preuves).
  \item Principe fondateur~: ne rien altérer. Bloqueur en écriture,
        copie sur image, hashes d'intégrité à chaque étape.
  \item Ordre de volatilité (RFC~3227)~: RAM~$\to$~réseau~$\to$~processus
        ~$\to$~disque~$\to$~journaux~$\to$~archives.
  \item Décision live/dead~: live pour capturer la RAM~; dead pour
        l'image disque intègre.
  \item Formats d'image~: RAW (dd), E01 (FTK/EnCase), AFF4.
        Outils~: dd, dcfldd, FTK Imager, avml/WinPmem.
  \item Chaîne de custody~: documenter chaque accès, chaque hash,
        chaque transfert.
  \item Environnement SIFT Workstation~: isolé réseau, lecture seule,
        journalisation des actions.
\end{itemize}
\end{boxresume}

\begin{boxexo}
\textbf{Ex.~5.1 — Stratégie d'acquisition.}
Vous arrivez devant \texttt{srv-erp-01} (Windows Server~2022) qui
est allumé et sur lequel une connexion RDP suspecte est active.
Décrivez dans l'ordre les 8~premières actions à effectuer, en
justifiant chaque décision.

\medskip\textbf{Ex.~5.2 — Commande d'acquisition.}
Rédigez la commande \texttt{dcfldd} complète pour acquérir le disque
\texttt{/dev/sdc} vers le fichier \texttt{erp-sdc.dd} sur un support
externe \texttt{/mnt/evidence}, avec~: hachage SHA-256, log du hash
dans \texttt{erp-sdc.sha256.log}, blocs de 128K.
Ajoutez la commande de vérification d'intégrité.

\medskip\textbf{Ex.~5.3 — PV de saisie.}
Rédigez le PV de collecte complet pour l'exercice~5.2,
en respectant le modèle de la boîte custody de ce chapitre.
\end{boxexo}
